% Baccalauréat malgache, série D, 2006 : énoncé.
\TCDBacAttribution{D}{2006}
\providecommand{\Courbe}{(\mathscr{C})}
\providecommand{\D}{(\mathscr{D})}
\providecommand{\Plan}{(\mathscr{P})}
\providecommand{\Ouv}{(O ; \vec{u} , \vec{v})}
%\providecommand{\Oij}{(O ; \vec{\imath} , \vec{\jmath})}
\providecommand{\cm}{~\mathrm{cm}}
\providecommand{\dx}{~\mathrm{d}x}
\providecommand{\Arg}[1]{\mathrm{Arg}\left( #1 \right)}
\providecommand{\Angl}[2]{\left( \overrightarrow{#1} ; \overrightarrow{#2} \right)}
%%%%%%%%%%%%%%%%%%%%%%%%%%%%%%%%%%%%%%%%%%%%%%%%%%%%%%%%%%%%%%%%%%%%%%%%%%%%%%%%%%%%%%%%%%%%%%%
%%%%%%%%%%%%%%%%%%%                   %%%%%%%%%%%%%%%%%%%                   %%%%%%%%%%%%%%%%%%%
%%%%%%%%%%%%%%%%%%%%%%%%%%%%%%%%%%%%%%%%%%%%%%%%%%%%%%%%%%%%%%%%%%%%%%%%%%%%%%%%%%%%%%%%%%%%%%%
\noindent \underline{\textsc{Exercice I}} (5 points)
      \begin{enumerate}
            \item R\'esoudre dans $\mathbb{C}$ l'\'equation : $z^2 - (1+i)\; z + 2 + i = 0$.
            \item Dans le plan complexe $\Plan$ muni d'un rep\`ere orthonorm\'e $\Ouv$, (unit\'e $1 \cm$), 
                  on donne les points $A$, $B$, $C$ et $D$ d'affixes respectives : 
                  $z_A = -1 + i$ ; $z_B = 1 - i$ ; $z_C = 2i$ et $z_D = 4 + 6i$.
                  \begin{enumerate}
                        \item Placer les points $A$, $B$, $C$ et $D$.
                        \item Ecrire l'expression complexe de la similitude directe $S$ telle que $S(D) = A$ et $S(C) = B$.
                        \item Donner les \'el\'ements caract\'eristiques de la similitude $S$.
                  \end{enumerate}
            \item \`A tout point $M$ d'affixe $z$ ($M \neq D$), on associe le point $M_1$ d'affixe $z_1$ tel que :
                  $z_1 = \frac{2 z - 4 i}{z - 4 - 6i}$.
                  \begin{enumerate}
                        \item Montrer que $OM_1 = \frac{2 CM}{DM}$.
                        \item D\'eterminer et repr\'esenter graphiquement les ensembles $(\Gamma)$ et $(\Gamma')$ tels que :
                              \begin{mitanisa}
                                    $(\Gamma)$ est l'ensemble des points $M$ d'affixe $z$ tels que $|z_1|=2$. \\
                                    $(\Gamma')$ est l'image de $(\Gamma)$ par la similitude $S$.
                              \end{mitanisa}
                  \end{enumerate}
      \end{enumerate}
%%%%%%%%%%%%%%%%%%%%%%%%%%%%%%%%%%%%%%%%%%%%%%%%%%%%%%%%%%%%%%%%%%%%%%%%%%%%%%%%%%%%%%%%%%%%%%%
%%%%%%%%%%%%%%%%%%%                   %%%%%%%%%%%%%%%%%%%                   %%%%%%%%%%%%%%%%%%%
%%%%%%%%%%%%%%%%%%%%%%%%%%%%%%%%%%%%%%%%%%%%%%%%%%%%%%%%%%%%%%%%%%%%%%%%%%%%%%%%%%%%%%%%%%%%%%%
\noindent \underline{\textsc{Exercice II}} (5 points) \newline
      Une urne contient 2 boules rouges et 4 boules vertes.
      Les boules rouges portent respectivement les num\'eros 0 et 1 ; 
      les boules vertes portent respectivement les num\'eros 2, 3, 4 et 5.
      L'\'epreuve $(E)$ consiste \`a tirer simultan\'ement 3 boules de l'urne ; 
      on appelle << succ\`es >> l'obtention de 3 boules de m\^eme couleur.
      \begin{enumerate}
            \item On effectue une fois l'\'epreuve $(E)$.
                  Calculer la probabilit\'e de chacun des \'ev\`enements suivants : 
                  \begin{itemize}
                        \item $A$ : << Avoir un succ\`es >>.
                        \item $B$ : << Avoir 3 boules dont la somme des num\'eros est \'egale \`a 6 >>.
                  \end{itemize}
            \item A chaque \'epreuve $(E)$, on d\'esigne par $X$ la variable al\'eatoire r\'eelle qui prend comme
                  valeur le plus grand des num\'eros port\'es sur les 3 boules tir\'ees.
                  \begin{enumerate}
                        \item Donner l'univers-image de $X$.
                        \item Donner la loi de probabilit\'e de $X$.
                  \end{enumerate}
            \item On r\'ep\`ete $n$ fois de suite, d'une mani\`ere ind\'ependante l'\'epreuve $(E)$.
                  On note l'\'ev\`enement : 
                  $A_n$ : << Obtenir au moins un succ\`es lors des $n$ \'epreuves $(E)$ >>.
                  \begin{enumerate}
                        \item Calculer la probabilit\'e $p_n$ de l'\'ev\`enement $A_n$.
                        \item D\'eterminer le nombre minimum d'\'epreuves $(E)$ qu'on doit effectuer pour que $p_n > 0,99$.
                  \end{enumerate}
                  %
                  On donne : $\ln \frac{5}{4} \simeq 0,22$ ; $\ln 10 \simeq 2,30$.
      \end{enumerate}
%%%%%%%%%%%%%%%%%%%%%%%%%%%%%%%%%%%%%%%%%%%%%%%%%%%%%%%%%%%%%%%%%%%%%%%%%%%%%%%%%%%%%%%%%%%%%%%
%%%%%%%%%%%%%%%%%%%                   %%%%%%%%%%%%%%%%%%%                   %%%%%%%%%%%%%%%%%%%
%%%%%%%%%%%%%%%%%%%%%%%%%%%%%%%%%%%%%%%%%%%%%%%%%%%%%%%%%%%%%%%%%%%%%%%%%%%%%%%%%%%%%%%%%%%%%%%
\noindent \underline{\textsc{Probleme}} (10 points) \newline
      On consid\`ere la fonction $f$ d\'efinie sur $]0;+\infty[$ par : 
      $f(x) = \frac{e}{x} + \frac{\ln x}{x^2}$ o\`u $e$ d\'esigne le nombre r\'eel base des logarithmes n\'ep\'eriens.
      On note $\Courbe$ sa courbe repr\'esentative dans un rep\`ere orthogonal direct \Oij{} tel que 
      $\| \ivec \| = 2 \cm $ et $\| \jvec \| = 1 \cm$.
      \begin{enumerate}
            \item On consid\`ere la fonction $g$ d\'efinie sur $]0;+\infty[$ par $g(x) = - 2 \ln x - x \; e + 1$.
                  \begin{enumerate}
                        \item Donner le sens de variation de $g$.
                        \item Montrer que l'\'equation $g(x) = 0$ admet une solution unique $\alpha$ dans l'intervalle
                              $\left] \frac{1}{2} ; 1 \right[$.
                        \item En d\'eduire que, pour tout $x \in ]0;\alpha[$, $g(x) > 0$ et pour tout 
                              $x \in ]\alpha ; +\infty[$, $g(x) < 0$.
                        \item Montrer que $f(\alpha) = \frac{1 + \alpha\; e}{2 \alpha^2}$.
                  \end{enumerate}
            \item \begin{enumerate}
                        \item Calculer les limites de $f$ en $0$ et en $+\infty$.
                        \item Calculer la d\'eriv\'ee $f'(x)$ de la fonction $f$ est v\'erifier que pour tout 
                              $x \in ]0 ; +\infty[$ : $f'(x) = \frac{g(x)}{x^3}$.
                        \item Dresser le tableau de variation de $f$.
                  \end{enumerate}
            \item Tracer la courbe $\Courbe$ en pr\'ecisant les branches infinies. \newline
                  On prendra : $\ln 2 = 0,69$ ; $e = 2,71$ ; $\alpha = 0,67$ et $f(\alpha) = 3,16$ pour la construction.
            \item Pour $n \in \mathbb{N}$, on pose 
                  $I_n = \int_{\expo{n}}^{\expo{n+1}} \frac{\ln x}{x^2} \dx$ et $I_n = \int_{\expo{n}}^{\expo{n+1}} f(x) \dx$.
                  \begin{enumerate}
                        \item A l'aide d'une int\'egration par parties, montrer que : 
                              $I_n = \frac{n+1}{\expo{n}} - \frac{n+2}{\expo{n+1}}$.
                        \item Montrer que $A_n = I_n + e$.
                        \item Calculer en $\cm^2$ l'aire du domaine plan compris entre la courbe $\Courbe$,
                              l'axe des abscisses et les droites d'\'equations respectives : $x = 1$ et $x = e$.
                  \end{enumerate}
      \end{enumerate}







